\section{Design Discussion: Training Run and Shortcut Properties}
\label{sec:appendix-run}

In ML-Schema, the relationship between a trained model and its algorithm is
mediated by \texttt{mls:Run}, which represents a training execution:
%
\begin{align*}
  \mathsf{mls{:}Algorithm}
    &\xleftarrow{\mathsf{mls{:}executes}}
    \mathsf{mls{:}Run}
    \xrightarrow{\mathsf{mls{:}hasOutput}}
    \mathsf{mls{:}Model} \\
  \mathsf{mls{:}Dataset}
    &\xleftarrow{\mathsf{mls{:}hasInput}}
    \mathsf{mls{:}Run}
\end{align*}
%
The \texttt{mls:Run} concept captures the full context of a training
execution---algorithm, input data, output model, hyperparameter settings,
and potentially runtime metadata (duration, hardware, random seeds).

The \MLIPsOntology{} follows this pattern with $\TrainingRunC$ (a subclass
of both \texttt{mls:Run} and \texttt{prov:Activity}), which mediates
between $\MLIPAlgorithmC$, $\TrainingDatasetC$, and $\TrainedModelC$ via
$\executesR$, $\trainedOnR$, and $\producesR$ respectively. Additionally,
we provide the shortcut roles $\trainedWithR$ and $\trainedOnR$ directly on
$\TrainedModelC$, defined as property chains through the run. These
shortcuts enable direct queries without navigating the run entity, which is
important because our primary data source---published MLIP
literature---rarely reports the operational details of training runs
(hardware, duration, random seeds). In many cases, the $\TrainingRunC$
entity will carry only its algorithm, dataset, and model links, making the
shortcut properties the more natural query path.

\paragraph{Critical assessment.}
The shortcut properties introduce redundancy: the same information is
accessible both directly on the model and via the run. While the property
chain axioms guarantee consistency in principle, in practice data entry
errors could create inconsistencies if shortcuts and run-mediated paths are
populated independently. We mitigate this through SHACL shapes that
validate consistency between the two paths.

An alternative design would omit the shortcuts entirely and always require
navigation through $\TrainingRunC$. This would be cleaner from an ontology
engineering perspective but would make common queries (``which algorithm
produced this model?'') unnecessarily verbose, particularly for data
extracted from publications where the run carries no additional information
beyond the links already captured by the shortcuts.
